\bookchapter{Capitolo 3: Collegare il telefono alla “linea di vita” — La rete}

\chapterintro{Senza rete, lo smartphone perde metà della sua “intelligenza”. La rete è la linea di vita del telefono moderno: capirla, usarla bene e proteggerla è una materia obbligatoria per ogni utente.}

\booksection{Non perdersi nella connessione — distinguere “strada” e “veicolo” tra Wi-Fi e traffico dati}

Molte persone vedono contemporaneamente accese le icone Wi-Fi e segnale sul telefono, ma sono totalmente confuse sulla differenza essenziale. L'equivoco più classico è:\textbf{sei connesso al Wi-Fi ma le pagine non si caricano}— perché il router ha perso la connessione, e il telefono non è tornato automaticamente al traffico dati. Per capire a fondo il problema, dobbiamo costruire un modello cognitivo “strada e veicolo”.

\begin{tcbinfo}{Metafora centrale: strada e veicolo}
Immagina Internet come una città. La\textbf{“strada”}è il canale tra il tuo telefono e Internet; il\textbf{“veicolo”}sono i dati che viaggiano sulla strada. Wi-Fi e traffico dati sono due strade diverse, ma i dati (i veicoli) sono gli stessi.

\end{tcbinfo}

\booksubsection{Differenza essenziale: canale locale vs autostrada remota}

Il Wi-Fi è essenzialmente una\textbf{tecnologia LAN}, come una strada senza fili sotto casa, che trasmette i dati dal router al telefono. Ma se la linea broadband non è pagata o il cavo ottico è rotto, la strada c'è ancora ma l'“autostrada” esterna è interrotta, e non si va da nessuna parte. Il traffico dati (dati mobili), invece, viaggia attraverso le stazioni base 4G/5G dell'operatore: è come salire direttamente sull'autostrada dedicata dell'operatore; finché c'è segnale, sei connesso a Internet.

\booksubsection{Costi e scenari: abbonamento fisso vs tariffazione a consumo}

Il Wi-Fi di solito è un investimento una tantum (canone mensile della linea): finché c'è corrente, tutta la famiglia lo usa senza limiti, ideale per scaricare film o giocare, scenari ad alto consumo di dati. Il traffico dati, invece, viene tariffato a consumo; superato il bundle può esserci rallentamento o addebiti extra, ed è più adatto a task leggeri in mobilità, come in metropolitana o in viaggio di lavoro.

\begin{tcbpractice}{Prova: controlla il tuo piano dati}
Apri [Impostazioni] > [Rete mobile] > [Gestione traffico] e verifica quanti GB ti restano questo mese e quali app consumano di più. Sapere cosa hai ti fa uscire di casa senza ansie.

\end{tcbpractice}

\booksubsection{Considerazioni di sicurezza: il Wi-Fi pubblico è un “campo di pesca” nascosto}

Per la sicurezza, i dati mobili sono in genere più sicuri del Wi-Fi pubblico. I dati mobili sono cifrati dall'operatore, mentre il Wi-Fi pubblico può essere facilmente trasformato in “hotspot falso” dai hacker per intercettare traffico non cifrato. Va però chiarito: oggi, con la diffusione capillare di HTTPS, l'attaccante può intercettare\textbf{quali siti hai visitato}(dominio) e i\textbf{contenuti HTTP non cifrati}; ma i messaggi WeChat, le password di login delle app bancarie e altri dati cifrati con HTTPS\textbf{non possono essere letti direttamente}— a meno che non venga messo in atto un attacco SSL stripping più complesso e l'utente ignori gli avvisi sui certificati. In ogni caso, per pagamenti fuori casa, usare direttamente i dati mobili resta la scelta più sicura.

\textbf{Promemoria di sicurezza}\par

In aeroporti, centri commerciali, caffè e altri luoghi pubblici, non connetterti a Wi-Fi gratuiti con nomi sospetti (es. “Free\_WiFi\_2024”). Anche se ti connetti, non usarli mai per accedere a un'app bancaria o per inserire password.

\booksubsection{Verifica lato telefono: capire il “linguaggio segreto” della barra di stato}

Come sapere cosa sta usando il telefono in questo momento? Su un telefono Huawei, ad esempio, la comparsa dell'\textbf{icona a doppia freccia}accanto al segnale dell'operatore o al segnale WLAN nella barra di stato indica che si sta usando quel metodo per la connessione; se entrambe sono accese, l'accelerazione di rete è attiva e il telefono sta usando entrambi i canali.

Per gli utenti iPhone è più diretto: se l'icona Wi-Fi è accesa, si sta usando il Wi-Fi; se è spenta, guarda le barre di segnale in alto a destra per capire che stai usando il traffico dati. Alcuni telefoni Android mostrano anche “5G” o “4G” nella barra di stato, indicando il tipo di rete mobile corrente.

\booksection{Reti mobili (4G/5G) e traffico dati}

La rete mobile è la base per connettersi ovunque. Da 2G a 5G la velocità è aumentata di migliaia di volte, ma il principio è simile.

\booksubsection{Standard di rete:}

\begin{itemize}
  \item \textbf{4G LTE}: attualmente la copertura più ampia, velocità sufficiente (teoricamente fino a 100 Mbps+), consumo energetico relativamente basso.
  \item \textbf{5G}: velocità elevata (teoricamente fino a 10 Gbps), ma la copertura è ancora in via di completamento, consumo energetico alto. Adatta a scaricare file grandi e guardare video ad alta risoluzione.
  \item \textbf{5G SA vs NSA}: SA è la rete autonoma, NSA si appoggia alle stazioni base 4G. L'esperienza SA è migliore, ma richiede il supporto sia dell'operatore sia del telefono.
\end{itemize}

\booksubsection{Gestione del traffico:}

\begin{itemize}
  \item \textbf{Cosa sono i dati}: è il volume di dati generato navigando sul telefono. Video e immagini consumano di più, testo e voce consumano meno.
  \item \textbf{Dritta per risparmiare}: attiva la "modalità risparmio dati" nelle impostazioni e aggiorna le app o guarda video ad alta risoluzione solo sotto Wi-Fi.
  \item \textbf{Barre di segnale $\neq$ velocità}: un segnale forte indica solo una buona connessione; la velocità dipende anche dal carico della stazione base e dalla tua distanza da essa.
\end{itemize}

\textbf{Avviso importante: prima di andare all'estero disattiva il [data roaming]!}\par

Usare i dati mobili all'estero può generare bollette astronomiche! Prima di partire, disattiva il data roaming nelle impostazioni oppure acquista in anticipo una scheda SIM locale o un pacchetto di roaming internazionale.

\booksection{Wi-Fi: connessione, condivisione e sicurezza}

Il Wi-Fi è la connessione principale in casa e in ufficio: veloce e senza consumo di traffico. Ma anche il Wi-Fi si evolve e le differenze tra gli standard sono notevoli.

\booksubsection{Connessione Wi-Fi:}

Vai in [Impostazioni] > [WLAN], scegli la rete e inserisci la password. Molti telefoni oggi supportano la connessione tramite scansione di codice QR, ancora più comoda.

\booksubsection{Condivisione della password:}

Sia Android che iOS permettono a un telefono già connesso di condividere rapidamente la password con un amico. Android genera un codice QR, iOS può condividere automaticamente per prossimità.

\booksubsection{Evoluzione degli standard Wi-Fi:}

Gli standard Wi-Fi sono definiti dall'IEEE; ogni generazione porta miglioramenti significativi:

\begin{center}
\begin{tabularx}{\linewidth}{|X|X|X|X|X|}
\hline
\textbf{Standard} & \textbf{Protocollo} & \textbf{Banda} & \textbf{Velocità massima} & \textbf{Caratteristiche} \tabularnewline
\hline
Wi-Fi 4 & 802.11n & 2.4/5GHz & 600Mbps & Standard datato, copertura discreta \tabularnewline
Wi-Fi 5 & 802.11ac & 5GHz & 3.5Gbps & Oggi il più diffuso, velocità adeguata \tabularnewline
Wi-Fi 6 & 802.11ax & 2.4/5GHz & 9.6Gbps & Prestazioni multi-dispositivo elevate, latenza bassa \tabularnewline
Wi-Fi 6E & 802.11ax & 2.4/5/6GHz & 9.6Gbps & Aggiunge la banda 6GHz, meno congestione \tabularnewline
Wi-Fi 7 & 802.11be & 2.4/5/6GHz & 46Gbps & Trasmissione ultra veloce, aggregazione multi-link \tabularnewline
\hline
\end{tabularx}
\end{center}

\begin{itemize}
  \item \textbf{Tecnologie chiave del Wi-Fi 6}: OFDMA (Orthogonal Frequency-Division Multiple Access) consente la trasmissione simultanea di più dispositivi; MU-MIMO (Multi-User Multiple Input Multiple Output) migliora l'efficienza multi-dispositivo; TWT (Target Wake Time) riduce i consumi.
  \item \textbf{Wi-Fi 6E}: aggiunge la banda 6GHz al Wi-Fi 6, offrendo più canali e meno interferenze. Richiede che telefono e router lo supportino entrambi.
  \item \textbf{Wi-Fi 7}: introduce MLO (Multi-Link Operation), può usare contemporaneamente più bande per trasmettere dati, riducendo drasticamente la latenza. Il canale ultra-largo a 320 MHz porta un salto nelle velocità. Piattaforme flagship come Snapdragon 8 Elite Gen5 e Dimensity 9400 supportano già il Wi-Fi 7; modelli come Xiaomi 17, Huawei Mate 80, Samsung S26 possono usare reti Wi-Fi 7 (in abbinamento a un router Wi-Fi 7).
  \item \textbf{2.4GHz vs 5GHz}: 2.4GHz attraversa meglio i muri ma è più lento e soggetto a interferenze; 5GHz è più veloce ma attraversa peggio i muri. La banda 6GHz del Wi-Fi 6E/7 è la più veloce, ma la più debole nell'attraversare i muri, adatta all'uso nella stessa stanza.
\end{itemize}

\booksubsection{Rischi del Wi-Fi pubblico:}

Il Wi-Fi pubblico può esporre a intercettazioni o a hotspot esca. Regole d'oro:

\begin{itemize}
  \item Evita di accedere a conti bancari o di inserire password.
  \item Usa preferibilmente una VPN per cifrare la connessione.
  \item Non connetterti automaticamente a Wi-Fi sconosciuti.
  \item Usa siti HTTPS (icona del lucchetto nella barra degli indirizzi).
\end{itemize}

\booksection{Bluetooth: l'onnipresente della comunicazione wireless a breve distanza}

Il Bluetooth è probabilmente una delle tecnologie wireless più usate sul tuo telefono — cuffie wireless, smartwatch, sistemi auto, tastiere e mouse si connettono tutti grazie a essa.

\booksubsection{Principio tecnico:}

Il Bluetooth opera nella banda ISM a 2.4GHz e usa la tecnica FHSS (Frequency Hopping Spread Spectrum) per distribuire i dati su 79 canali a salti rapidi, evitando le interferenze. La distanza di comunicazione è divisa in tre classi: Classe 1 circa 100 m, Classe 2 circa 10 m (comune sui telefoni), Classe 3 circa 1 m.

\booksubsection{Evoluzione delle versioni:}

\begin{itemize}
  \item \textbf{Bluetooth 4.0}: introduce il Bluetooth a basso consumo (BLE), consentendo a dispositivi con batteria piccola come braccialetti e fasce cardiache di connettersi senza fili; l'autonomia passa da giorni a mesi.
  \item \textbf{Bluetooth 4.2}: migliora velocità di trasmissione e protezione della privacy, supporta la connessione diretta IPv6.
  \item \textbf{Bluetooth 5.0}: velocità raddoppiata (2 Mbps), portata quadrupla (teoricamente 300 m), capacità di broadcast otto volte superiore. Supporta il doppio audio (due cuffie contemporanee).
  \item \textbf{Bluetooth 5.1}: introduce la funzione di indirizzamento direzionale (AoA/AoD), con precisione di localizzazione centimetrica.
  \item \textbf{Bluetooth 5.2}: nuovo standard LE Audio, supporta la codifica LC3 (qualità audio migliore e consumi più bassi), audio multi-flusso e audio broadcast.
  \item \textbf{Bluetooth 5.3/5.4}: ottimizzazioni ulteriori su consumi e stabilità della connessione, broadcast periodico migliorato.
\end{itemize}

\textbf{Diffusione del LE Audio}: nel 2025-2026, il LE Audio è entrato in una fase di rapida diffusione. Le nuove cuffie true wireless dei principali brand (come Huawei FreeBuds Pro 5, Xiaomi Buds 6 Pro, Samsung Galaxy Buds3 Pro) supportano già LE Audio e la codifica LC3, offrendo latenza più bassa e qualità audio superiore. Quando acquisti nuove cuffie, è consigliabile scegliere prodotti che supportano LE Audio.

\booksubsection{Scenari applicativi:}

\begin{itemize}
  \item \textbf{Trasmissione audio}: cuffie wireless, impianti auto. Il LE Audio del Bluetooth 5.2 porta un salto qualitativo in audio e latenza.
  \item \textbf{Trasmissione dati}: scambio di file e foto a breve distanza. Pur essendo più lento del Wi-Fi Direct, vince per semplicità.
  \item \textbf{Accoppiamento dispositivi}: smartwatch, braccialetti, tastiere e mouse, gamepad.
  \item \textbf{Servizi di posizione}: Beacon Bluetooth per navigazione indoor e tracciamento oggetti (es. Apple AirTag).
  \item \textbf{Casa intelligente}: rete Bluetooth Mesh per controllare lampadine, serrature ecc.
\end{itemize}

\booksection{Telecommando a infrarossi: il telefono diventa un telecomando universale}

Il telecommando a infrarossi è una delle tecnologie di controllo wireless più antiche, ma ancora utile sui telefoni. Alcuni telefoni Android (come alcuni modelli Xiaomi e Huawei) integrano un trasmettitore IR, in grado di controllare TV, condizionatori, decoder e altri elettrodomestici.

\booksubsection{Principio di funzionamento:}

Il telecommando IR emette, tramite un LED a infrarossi, impulsi luminosi a una frequenza specifica (in genere 38 kHz); il ricevitore IR dell'elettrodomestico decodifica gli impulsi ed esegue l'azione corrispondente. Marche diverse usano protocolli di codifica diversi, tra i più comuni NEC, RC5, RC6, SONY. L'app IR del telefono integra il database di codici di ciascun marchio: basta selezionare marca e modello per la corrispondenza.

\booksubsection{Limiti:}

\begin{itemize}
  \item \textbf{Direzionalità}: gli infrarossi viaggiano in linea retta, occorre puntare il dispositivo, non attraversano i muri.
  \item \textbf{Breve distanza}: portata utile in genere 5-8 metri.
  \item \textbf{Non supportato da Apple}: l'iPhone non ha il trasmettitore IR e non può usare questa funzione.
\end{itemize}

\booksection{NFC: la magia del “tocca e vai”}

NFC (Near Field Communication) è una tecnologia di comunicazione wireless ad alta frequenza a brevissima distanza (entro 10 cm), operante a 13.56 MHz, con velocità massima di 424 Kbps.

\booksubsection{Principio di comunicazione e modalità:}

NFC deriva dalla tecnologia RFID e supporta tre modalità di funzionamento:

\begin{itemize}
  \item \textbf{Modalità lettura/scrittura}: il telefono legge informazioni da un tag NFC, ad esempio un adesivo NFC che attiva un'azione.
  \item \textbf{Modalità peer-to-peer}: due dispositivi NFC si scambiano dati, come Android Beam (discontinuato).
  \item \textbf{Modalità card emulation}: il telefono emula una smart card — è la modalità più usata: carte bus, badge di accesso, carte bancarie si basano su di essa.
\end{itemize}

\booksubsection{Meccanismi di sicurezza:}

\begin{itemize}
  \item \textbf{Brevissima distanza}: la comunicazione entro 10 cm è di per sé una garanzia di sicurezza fisica; l'intercettazione remota è praticamente impossibile.
  \item \textbf{Chip di sicurezza}: il telefono integra un chip di sicurezza dedicato (SE) o usa HCE (Host Card Emulation); le chiavi di pagamento sono conservate in un'area sicura e non possono essere estratte nemmeno con il telefono in root.
  \item \textbf{Verifica impronta/volto}: il pagamento richiede conferma biometrica, così anche se il telefono viene smarrito non si può frodare.
  \item \textbf{Tokenizzazione}: Apple Pay, Huawei Pay e simili non trasmettono il numero reale della carta, ma un token monouso; anche se intercettato, non può essere riusato.
\end{itemize}

\booksubsection{Applicazioni tipiche:}

\begin{itemize}
  \item \textbf{Carta bus / badge di accesso}: si scrivono nel telefono le informazioni della carta fisica e si usa il telefono per transitare.
  \item \textbf{Pagamento mobile}: Apple Pay, Huawei Pay ecc., collegando la carta bancaria al telefono per pagare con un tocco.
  \item \textbf{Accoppiamento dispositivi}: NFC per accoppiare rapidamente cuffie e casse Bluetooth.
  \item \textbf{Autenticazione d'identità}: lettura di carte d'identità NFC, verifica della patente elettronica.
\end{itemize}

\booksection{Hotspot personale: trasforma il telefono in un router Wi-Fi}

Quando hai molto traffico dati, puoi attivare l'hotspot per tablet e computer.

\begin{itemize}
  \item \textbf{Come attivarlo}: Impostazioni > Hotspot personale / Tethering, e attivalo.
  \item \textbf{Note}: consuma molta batteria e traffico; imposta una password complessa e limita il numero di dispositivi connessi.
  \item \textbf{Tethering USB}: collega il computer con il cavo dati, più stabile e veloce del hotspot wireless.
\end{itemize}

\begin{tcbpractice}{Prova sul campo: configura e gestisci l'hotspot}
\begin{enumerate}
  \item Apri [Impostazioni], trova "Hotspot personale" o "Tethering", attivalo e imposta nome e password.
  \item Con un altro dispositivo (ad esempio il computer) cerca e connettiti a questo Wi-Fi.
  \item Sul telefono verifica se riesci a gestire i dispositivi connessi o a espellere quelli sconosciuti.
\end{enumerate}

\end{tcbpractice}
